% 03-recording.tex

\section{Recording and Representation}\label{sec:recording}

The recording pipeline, graph construction, and Merkle DAG are implemented and tested.
Formal properties of the DAG are proved in Lean~4.
Figure~\ref{fig:pipeline} shows the end-to-end data flow.

\begin{figure}[ht]
\centering
\begin{tikzpicture}[
  node distance=1.4em,
  box/.style={draw=black!40, fill=black!3, rounded corners=1.5pt,
    minimum height=2.2em, font=\sffamily\scriptsize, align=center},
  buf/.style={box, fill=blue!5},
  obox/.style={box, fill=orange!5},
  arr/.style={-{Stealth[length=4pt]}, thick, black!40},
  thr/.style={font=\sffamily\tiny, text=black!40},
]
  %% Foreground
  \node[box, fill=blue!6] (vessel) {Vessel\\(ATen dispatch)};
  \node[buf, right=1.8em of vessel] (ring) {TraceRing\\64\,B $\times$ $2^{16}$};
  \node[buf, below=0.7em of ring]   (meta) {MetaLog\\144\,B $\times$ $2^{20}$};

  %% Background
  \node[box, fill=green!5, right=2.8em of ring] (bg) {Background\\Thread};
  \node[box, right=1.6em of bg]   (graph) {TraceGraph\\(bidir.\ CSR)};
  \node[box, right=1.6em of graph] (dag) {Merkle\\DAG};

  %% Outputs
  \node[obox, below=0.8em of dag]   (cache) {KernelCache};
  \node[obox, below=0.8em of graph] (plan)  {MemoryPlan};

  %% Arrows
  \draw[arr] (vessel) -- (ring);
  \draw[arr] (vessel) |- (meta);
  \draw[arr] (ring)   -- (bg);
  \draw[arr] (meta)   -| (bg);
  \draw[arr] (bg)     -- (graph);
  \draw[arr] (graph)  -- (dag);
  \draw[arr] (dag)    -- (cache);
  \draw[arr] (graph)  -- (plan);

  %% Thread boundary
  \draw[dashed, black!15, thick]
    ($(ring.north east)!0.5!(bg.north west)+(0,1.2)$) --
    ($(meta.south east)!0.5!(bg.south west)+(0,-0.7)$);
  \node[thr] at ($(vessel.north)+(0,0.9)$) {foreground};
  \node[thr] at ($(bg.north)+(0,0.9)$) {background};

  %% Timing
  \node[font=\tiny, text=black!40, below=0.2em of vessel] {${\sim}5$\,ns/op};
\end{tikzpicture}
\caption{Recording pipeline.
  The foreground thread writes to two lock-free SPSC buffers.
  The background thread drains, builds a bidirectional CSR graph,
  constructs the Merkle DAG, and computes a static memory plan.}
\label{fig:pipeline}
\end{figure}

\subsection{Vessel Interception}

The Vessel adapts Crucible to the host framework.
The current Vessel targets PyTorch's ATen dispatch mechanism---a priority-ordered function pointer table indexed by dispatch keys.
Crucible registers \code{DispatchKey::Conductor} above backend keys, observing every operation before it reaches the compute backend.

\textbf{Recording mode} (six steps per operation):
(1)~snapshot input \code{TensorMeta} (\code{sizes[8]}, \code{strides[8]}, \code{data\_ptr}, \code{ndim}, \code{dtype}, \code{device\_type}, \code{device\_idx}, \code{layout}---144 bytes per tensor);
(2)~compute \code{SchemaHash} (from registered operator name) and \code{ShapeHash} (from input dimensions);
(3)~execute the operation via redispatch;
(4)~snapshot output \code{TensorMeta};
(5)~append metadata to the MetaLog;
(6)~record a 64-byte entry in the TraceRing.
\textbf{Compiled mode}: advance operation index, check guard hash, return pre-allocated shadow handle.
No execution, no allocation, no redispatch.

The Vessel handles operations with zero tensor arguments (profiler hooks like \code{profiler::\_record\_\allowbreak{}function\_enter\_new}, scope markers) by recording them with \code{MetaIndex::none()}.
TensorList inputs (concatenation, stacking) are unpacked into individual \code{TensorMeta} entries.
Scalar arguments are encoded inline (up to five \code{int64} values per TraceRing entry, covering 99.9\% of operations).
Communication operations (all-reduce, all-gather) use the identical recording pipeline.

\subsection{Trace Recording}

Two parallel SPSC (single-producer, single-consumer) structures communicate between the foreground thread (Python + Vessel recording) and the background thread (graph construction + compilation).

\textbf{TraceRing.}
65,536 entries (\code{CAPACITY = 1 << 16}), each exactly 64 bytes (one cache line): \code{SchemaHash} (8B), \code{ShapeHash} (8B), \code{num\_inputs} (2B), \code{num\_outputs} (2B), \code{num\_scalar\_args} (2B), \code{grad\_enabled} (1B), \code{inference\_mode} (1B), and \code{scalar\_values[5]} (40B).
Three parallel arrays alongside entries---\code{meta\_starts[]} (\code{MetaIndex}), \code{scope\_hashes[]} (\code{ScopeHash}), \code{callsite\_hashes[]} (\code{CallsiteHash})---are indexed by the same slot, providing the background thread with module context and source location.
Total pre-allocated: ${\sim}5.25$\,MB.

Bitmask indexing (\code{index \& MASK}) replaces modulo.
Head and tail are \code{alignas(64)} atomics on separate cache lines to prevent false sharing.
The producer caches the last-read tail locally (\code{cached\_tail\_}); since drain only advances the tail forward, a stale cache is conservative.
This avoids cross-core cache-line traffic on the common path---approximately 20,000 appends between drains at 5\,ns/op and 100\,\textmu{}s drain intervals.
The next write slot is prefetched via \code{\_\_builtin\_prefetch} (write hint, L1 locality) before publishing the head, hiding DRAM latency for the 4\,MB ring.

\textbf{MetaLog.}
1,048,576 entries (\code{CAPACITY = 1 << 20}) of 144-byte \code{TensorMeta} records, 64-byte aligned allocation (${\sim}144$\,MB).
Same SPSC protocol with cached tail.
Producer hot path: ${\sim}12$\,ns for 1 tensor, ${\sim}33$\,ns for 3 tensors (measured).
Bulk memcpy for contiguous regions compiles to AVX-512 stores.
The background thread accesses entries zero-copy via \code{try\_contiguous()} when the range does not wrap around the circular boundary (99.99\% of calls with 1M capacity and ${\sim}1{,}500$ metas per iteration).

Both structures are pre-allocated at initialization and never resized.
The foreground recording path involves no dynamic allocation, no system calls, and no contention.
All graph and DAG memory is allocated from a bump-pointer Arena on the background thread (${\sim}2$\,ns per allocation, 1\,MB blocks, zero fragmentation).

\subsection{Graph Construction}

The background thread drains the TraceRing in batches of up to 4,096 entries (\code{BATCH\_SIZE}) and constructs a \code{TraceGraph}: a bidirectional CSR property graph where nodes represent operations and edges represent relationships.
Four edge kinds are defined: \code{DATA\_FLOW} (producer-consumer), \code{ALIAS} (shared storage), \code{CONTROL\_FLOW} (explicit ordering), and \code{SCALAR\_FLOW} (scalar dependencies).
Each Edge is 12 bytes: \code{OpIndex src}, \code{OpIndex dst}, \code{src\_port} (\code{uint8}), \code{dst\_port} (\code{uint8}), \code{EdgeKind}, padding.

\textbf{Data flow edges} are constructed by pointer tracking.
A PtrMap (open-addressing, generation-counter-based to avoid memset between iterations) maps tensor \code{data\_ptr} to the \code{OpIndex} that produced it.
Input pointer matching a recorded output pointer creates a \code{DATA\_FLOW} edge.
\textbf{Alias edges} arise when multiple operations reference the same underlying storage (views, in-place ops, transposes).

Each graph node is a \code{TraceEntry}: \code{SchemaHash}, \code{ShapeHash}, \code{ScopeHash}, \code{CallsiteHash} (all strong \code{uint64} types), arena-allocated \code{TensorMeta} arrays for inputs and outputs, arena-allocated scalar args, grad/inference flags, \code{CKernelId} classification, and \code{SlotId} arrays for memory planning.
The graph is built in a single pass with $O(V{+}E)$ CSR construction via counting sort (prefix sum + scatter).
Both forward (src-sorted) and reverse (dst-sorted) CSR adjacency structures are constructed simultaneously.

Scratch buffers (PtrMap, SlotInfo, Edge arrays) are allocated once and reused across iterations---zero per-call allocation.
Content hash is fused into the main loop as a streaming accumulator, eliminating a redundant second pass.

\textbf{Iteration detection.}
The \code{IterationDetector} uses a $K{=}5$ schema hash signature occupying exactly two cache lines (128 bytes).
Steady-state hot path: one increment + one comparison + one well-predicted branch (${\sim}1$\,ns).
The expected next hash value is cached directly---no rolling fingerprint, no ring buffer, no multiply.
Two-match confirmation handles warmup: the first match locks the signature as a candidate; the second confirms a boundary.
The last completed iteration length is computed via \code{std::sub\_sat(ops\_since\_boundary, K)}.
False positives require 5 consecutive hash collisions; with 64-bit hashes, probability is ${\sim}5 \times 2^{-320}$.

\subsection{The Merkle DAG}

From the \code{TraceGraph}, the background thread constructs the Merkle DAG.
All nodes inherit from \code{TraceNode} (24 bytes: \code{TraceNodeKind kind}, padding, \code{MerkleHash merkle\_hash}, \code{TraceNode* next}).
Four kinds: \code{REGION}, \code{BRANCH}, \code{LOOP}, \code{TERMINAL}.

\textbf{RegionNode.}
A compilable operation sequence.
Contains an arena-allocated array of \code{TraceEntry} structs, a \code{ContentHash} computed via \code{compute\_content\_hash()} (wymix-fold over schema hashes, per-tensor dimension products, packed dtype/device metadata, and scalar arguments), an atomic \code{CompiledKernel*} pointer (background writes with release, foreground reads with acquire), a \code{MemoryPlan*} for liveness analysis, \code{first\_op\_schema} for quick mismatch detection, and \code{measured\_ms}/\code{variant\_id} for runtime profiling.
The Merkle hash extends the content hash with the continuation's Merkle hash via \code{fmix64(content\_hash \^{} next->merkle\_hash)}.

\textbf{BranchNode.}
A guard point with sorted arms for $O(\log n)$ binary search during replay.
The \code{Guard} struct (12 bytes) supports five kinds: \code{SHAPE\_DIM}, \code{SCALAR\_VALUE}, \code{DTYPE}, \code{DEVICE}, \code{OP\_SEQUENCE}.
Each arm is a \code{(int64\_t value, TraceNode* target)} pair.
The \code{next} pointer points to the merge point where all arms reconverge.
BranchNodes are the mechanism for all runtime adaptation: architecture mutation, attention head replacement, hyperparameter changes, continuous learning with A/B verification.

\textbf{LoopNode.}
Cyclic computation within the acyclic DAG framework (64 bytes = one cache line).
Contains a body sub-DAG (linked list ending in \code{TERMINAL}), an array of \code{FeedbackEdge} structs (4 bytes each: \code{output\_idx}, \code{input\_idx}), and termination: \code{LoopTermKind::REPEAT} (fixed $N$ iterations) or \code{LoopTermKind::UNTIL} (converge within epsilon).
The Merkle hash uses a ``LOOPNODE'' salt (\code{0x4C4F4F504E4F4445}) to distinguish from RegionNodes, incorporating \code{body\_content\_hash}, \code{feedback\_signature()} (\code{fmix64} fold over all feedback edges with edge count), and \code{loopterm\_hash()} (termination kind + repeat count + epsilon bits).
LoopNodes enable compiled recurrence (RNNs without Python per timestep), convergence-based execution (fixed-point iterations), and cross-iteration pipelining.

\textbf{Atomic activation.}
The background thread prepares new DAG structures and publishes them via \code{active\_region.store(ptr, release)}.
The foreground thread loads with acquire.
All zero-downtime updates---compilation activation, branch swaps, memory plan updates, topology changes, rollbacks---use this single mechanism.

\begin{figure}[ht]
\centering
\begin{tikzpicture}[
  node distance=1.6em,
  R/.style={draw=blue!40, fill=blue!5, rounded corners=1pt,
    minimum width=5.5em, minimum height=1.8em, font=\sffamily\scriptsize},
  B/.style={draw=orange!50, fill=orange!6, rounded corners=1pt,
    minimum width=5.5em, minimum height=1.8em, font=\sffamily\scriptsize},
  L/.style={draw=green!50!black, fill=green!6, rounded corners=1pt,
    minimum width=5.5em, minimum height=1.8em, font=\sffamily\scriptsize},
  arr/.style={-{Stealth[length=3.5pt]}, black!35},
]
  \node[R] (r1) {RegionNode};
  \node[B, below=1.6em of r1] (b1) {BranchNode};
  \node[R, below left=1.4em and 2.5em of b1]  (a1) {Arm\,1};
  \node[R, below=1.4em of b1]                 (a2) {Arm\,2};
  \node[R, below right=1.4em and 2.5em of b1] (a3) {Arm\,3};
  \node[R, below=2.2em of a2] (merge) {RegionNode};
  \node[L, below=1.6em of merge] (loop) {LoopNode};

  \draw[arr] (r1) -- (b1);
  \draw[arr] (b1) -- (a1);
  \draw[arr] (b1) -- (a2);
  \draw[arr] (b1) -- (a3);
  \draw[arr] (a1) |- (merge);
  \draw[arr] (a2) -- (merge);
  \draw[arr] (a3) |- (merge);
  \draw[arr] (merge) -- (loop);
  \draw[arr, dashed, bend right=55] (loop.west)
    to node[left, font=\scriptsize\itshape, text=black!40] {feedback} (loop.north west);

  %% Annotations
  \node[font=\scriptsize, text=black!40, right=0.4em of b1]
    {guard: \texttt{SHAPE\_DIM}};
  \node[font=\scriptsize, text=black!40, right=0.4em of loop]
    {\texttt{REPEAT}($N$) $|$ \texttt{UNTIL}($\varepsilon$)};
\end{tikzpicture}
\caption{Merkle DAG node types.
  \textsf{RegionNode}: compilable op sequence.
  \textsf{BranchNode}: guard dispatch with sorted arms.
  \textsf{LoopNode}: cyclic computation with feedback edges.
  All modifications---architecture mutation, precision changes,
  continuous learning---are BranchNode arms.}
\label{fig:dagstructure}
\end{figure}

\subsection{Content Addressing and the Kernel Cache}

The \code{KernelCache} is a lock-free open-addressing hash table (default capacity 4,096, power-of-two, calloc-initialized).
Lookup: linear probing with bitmask indexing; \code{content\_hash} slots are \code{atomic<uint64\_t>} loaded with acquire; zero signals empty (miss).
Insert: CAS on the \code{content\_hash} slot with acq\_rel ordering; if the key already exists, the kernel pointer is updated (newer variant wins).
The size counter uses relaxed ordering (informational only).

Content addressing yields four reuse levels:

\begin{itemize}
  \item \textbf{Cross-iteration.} Iteration $N{+}1$ typically has identical content hashes. All compiled kernels are cache hits.
  \item \textbf{Cross-run.} The KernelCache persists in the Cipher. Restarting training incurs zero compilation.
  \item \textbf{Cross-Vigil.} Different Vigils sharing sub-computations (e.g., same attention mechanism, same head dimension) produce identical content hashes. Compiled kernels transfer automatically.
  \item \textbf{Cross-device.} The same content hash maps to different compiled kernels on different hardware. Multiple device-specific variants coexist.
\end{itemize}

The cache grows monotonically and persists across Relay migrations.
Figure~\ref{fig:content-addressing} illustrates the full content-addressing flow.

\begin{figure}[ht]
\centering
\begin{tikzpicture}[
  node distance=1em,
  box/.style={draw=black!35, fill=white, rounded corners=2pt,
    minimum height=2.4em, font=\sffamily\scriptsize, align=center,
    blur shadow={shadow blur steps=4, shadow xshift=0.4pt,
                 shadow yshift=-0.4pt, shadow blur radius=1.2pt}},
  hash/.style={box, fill=blue!6, minimum width=6.5em},
  cache/.style={box, fill=orange!6, minimum width=9em, minimum height=3.5em},
  inp/.style={box, fill=black!4, minimum width=7em, text width=6em},
  arr/.style={-{Stealth[length=5pt]}, thick, black!35},
  lbl/.style={font=\sffamily\tiny, text=black!45},
]
  %% Inputs
  \node[inp] (ops) {schema hashes\\shapes, strides\\dtypes, devices\\scalar args};

  %% Hashing
  \node[hash, right=2.5em of ops] (ch) {\textbf{ContentHash}\\$\in\mathcal{B}_{64}$};
  \node[hash, below=1.2em of ch]  (mh) {\textbf{MerkleHash}\\$\in\mathcal{B}_{64}$};

  %% KernelCache
  \node[cache, right=3em of ch] (kc) {\textbf{KernelCache}\\[1pt]
    {\tiny\texttt{(ContentHash, device\_cap)}}\\[-1pt]
    {\tiny$\longrightarrow$ \texttt{CompiledKernel*}}};

  %% Arrows
  \draw[arr] (ops) -- node[above, lbl] {wymix fold} (ch);
  \draw[arr] (ch)  -- node[above, lbl] {lock-free lookup} (kc);
  \draw[arr] (ch)  -- node[right, lbl, xshift=2pt]
    {$\oplus$ child hashes} (mh);

  %% Reuse levels (right of cache)
  \node[right=1.5em of kc, font=\sffamily\scriptsize, text=black!55,
        align=left, text width=8em] (reuse) {%
    \raisebox{0pt}{\tikz\fill[blue!40] (0,0) circle (2pt);}\;
      cross-iteration\\
    \raisebox{0pt}{\tikz\fill[green!50!black] (0,0) circle (2pt);}\;
      cross-run\\
    \raisebox{0pt}{\tikz\fill[orange!60] (0,0) circle (2pt);}\;
      cross-model\\
    \raisebox{0pt}{\tikz\fill[red!40] (0,0) circle (2pt);}\;
      cross-device};

  %% Structural diff annotation
  \node[lbl, below=0.3em of mh, text width=9em, align=center]
    {structural diff in $O(\log N)$\\--- shared subtrees identical};

  %% Brace for inputs
  \draw[decorate, decoration={brace, mirror, amplitude=4pt},
        thick, black!25]
    ([yshift=-2pt]ops.south west) -- ([yshift=-2pt]ops.south east)
    node[midway, below=5pt, lbl] {per-RegionNode};
\end{tikzpicture}
\caption{Content addressing.
  A RegionNode's semantics are folded into a 64-bit ContentHash via wymix.
  The KernelCache maps \texttt{(ContentHash, device\_capability)} to compiled
  artifacts, enabling four levels of reuse.
  The MerkleHash extends to the full subtree for $O(\log N)$ structural diff.}
\label{fig:content-addressing}
\end{figure}
